%% ============================================================
%% ebook.latex — Pandoc LaTeX template for Code4Projects ebooks
%% ============================================================
\documentclass[11pt,american,a4paper,]{book}

%% ── Encoding & language ─────────────────────────────────────
\usepackage[utf8]{inputenc}
\usepackage[T1]{fontenc}
\usepackage[shorthands=off,main=american]{babel}

%% ── Fonts (xelatex) ─────────────────────────────────────────
\usepackage{fontspec}
\setmainfont{Georgia}
\setsansfont{Helvetica Neue}
\setmonofont[Scale=0.88]{Menlo}

%% ── Page geometry ───────────────────────────────────────────
\usepackage[
  left=2.5cm,
  right=2.5cm,
  top=3cm,
  bottom=3cm,
]{geometry}

%% ── Line spacing ────────────────────────────────────────────
\usepackage{setspace}
\setstretch{1.4}

%% ── Colors & links ──────────────────────────────────────────
\usepackage[dvipsnames,svgnames,x11names]{xcolor}
\usepackage[
  colorlinks=true,
  linkcolor=NavyBlue,
  urlcolor=NavyBlue,
  citecolor=NavyBlue,
  filecolor=NavyBlue,
  pdfusetitle
]{hyperref}

%% ── Code listings ───────────────────────────────────────────
\usepackage{fancyvrb}
\usepackage{fvextra}

\definecolor{codebg}{HTML}{1e1e2e}
\definecolor{codefg}{HTML}{cdd6f4}
\definecolor{codeframe}{HTML}{313244}

\RecustomVerbatimEnvironment{verbatim}{Verbatim}{
  fontsize=\small,
  frame=single,
  framesep=6pt,
  rulecolor=\color{codeframe},
  fillcolor=\color{codebg},
  formatcom=\color{codefg},
  breaklines=true,
  breakanywhere=true,
  xleftmargin=10pt,
  xrightmargin=10pt,
}

%% ── Tables ──────────────────────────────────────────────────
\usepackage{longtable}
\usepackage{booktabs}
\usepackage{array}

%% ── Images ──────────────────────────────────────────────────
\usepackage{graphicx}
\makeatletter
\def\maxwidth{\ifdim\Gin@nat@width>\linewidth\linewidth\else\Gin@nat@width\fi}
\def\maxheight{\ifdim\Gin@nat@height>\textheight\textheight\else\Gin@nat@height\fi}
\makeatother
\setkeys{Gin}{width=\maxwidth,height=\maxheight,keepaspectratio}

%% ── Misc ────────────────────────────────────────────────────
\usepackage{parskip}
\setlength{\parindent}{0pt}
\usepackage{microtype}
\usepackage{enumitem}
\usepackage{titlesec}
\usepackage{mdframed}

%% Chapter style
\titleformat{\chapter}[display]
  {\normalfont\huge\bfseries\sffamily}
  {\chaptertitlename\ \thechapter}{16pt}{\Huge}
\titlespacing*{\chapter}{0pt}{30pt}{20pt}

%% Section style
\titleformat{\section}
  {\normalfont\Large\bfseries\sffamily}
  {\thesection}{1em}{}
\titleformat{\subsection}
  {\normalfont\large\bfseries\sffamily}
  {\thesubsection}{1em}{}

%% ── Headers / footers ───────────────────────────────────────
\usepackage{fancyhdr}
\pagestyle{fancy}
\fancyhf{}
\fancyhead[LE,RO]{\small\thepage}
\fancyhead[RE]{\small\itshape Getting Started with Docker}
\fancyhead[LO]{\small\itshape\leftmark}
\renewcommand{\headrulewidth}{0.4pt}

%% ── Blockquotes ─────────────────────────────────────────────
\newmdenv[
  leftline=true, rightline=false, topline=false, bottomline=false,
  linewidth=3pt, linecolor=NavyBlue!60,
  backgroundcolor=NavyBlue!5,
  innerleftmargin=12pt, innerrightmargin=12pt,
  innertopmargin=8pt, innerbottommargin=8pt,
  skipabove=10pt, skipbelow=10pt,
]{quoteblock}
\renewenvironment{quote}{\begin{quoteblock}}{\end{quoteblock}}

%% ── pandoc required macros ──────────────────────────────────
\providecommand{\tightlist}{%
  \setlength{\itemsep}{2pt}\setlength{\parskip}{0pt}}

\usepackage{color}
\usepackage{fancyvrb}
\newcommand{\VerbBar}{|}
\newcommand{\VERB}{\Verb[commandchars=\\\{\}]}
\DefineVerbatimEnvironment{Highlighting}{Verbatim}{commandchars=\\\{\}}
% Add ',fontsize=\small' for more characters per line
\newenvironment{Shaded}{}{}
\newcommand{\AlertTok}[1]{\textcolor[rgb]{1.00,0.00,0.00}{\textbf{#1}}}
\newcommand{\AnnotationTok}[1]{\textcolor[rgb]{0.38,0.63,0.69}{\textbf{\textit{#1}}}}
\newcommand{\AttributeTok}[1]{\textcolor[rgb]{0.49,0.56,0.16}{#1}}
\newcommand{\BaseNTok}[1]{\textcolor[rgb]{0.25,0.63,0.44}{#1}}
\newcommand{\BuiltInTok}[1]{\textcolor[rgb]{0.00,0.50,0.00}{#1}}
\newcommand{\CharTok}[1]{\textcolor[rgb]{0.25,0.44,0.63}{#1}}
\newcommand{\CommentTok}[1]{\textcolor[rgb]{0.38,0.63,0.69}{\textit{#1}}}
\newcommand{\CommentVarTok}[1]{\textcolor[rgb]{0.38,0.63,0.69}{\textbf{\textit{#1}}}}
\newcommand{\ConstantTok}[1]{\textcolor[rgb]{0.53,0.00,0.00}{#1}}
\newcommand{\ControlFlowTok}[1]{\textcolor[rgb]{0.00,0.44,0.13}{\textbf{#1}}}
\newcommand{\DataTypeTok}[1]{\textcolor[rgb]{0.56,0.13,0.00}{#1}}
\newcommand{\DecValTok}[1]{\textcolor[rgb]{0.25,0.63,0.44}{#1}}
\newcommand{\DocumentationTok}[1]{\textcolor[rgb]{0.73,0.13,0.13}{\textit{#1}}}
\newcommand{\ErrorTok}[1]{\textcolor[rgb]{1.00,0.00,0.00}{\textbf{#1}}}
\newcommand{\ExtensionTok}[1]{#1}
\newcommand{\FloatTok}[1]{\textcolor[rgb]{0.25,0.63,0.44}{#1}}
\newcommand{\FunctionTok}[1]{\textcolor[rgb]{0.02,0.16,0.49}{#1}}
\newcommand{\ImportTok}[1]{\textcolor[rgb]{0.00,0.50,0.00}{\textbf{#1}}}
\newcommand{\InformationTok}[1]{\textcolor[rgb]{0.38,0.63,0.69}{\textbf{\textit{#1}}}}
\newcommand{\KeywordTok}[1]{\textcolor[rgb]{0.00,0.44,0.13}{\textbf{#1}}}
\newcommand{\NormalTok}[1]{#1}
\newcommand{\OperatorTok}[1]{\textcolor[rgb]{0.40,0.40,0.40}{#1}}
\newcommand{\OtherTok}[1]{\textcolor[rgb]{0.00,0.44,0.13}{#1}}
\newcommand{\PreprocessorTok}[1]{\textcolor[rgb]{0.74,0.48,0.00}{#1}}
\newcommand{\RegionMarkerTok}[1]{#1}
\newcommand{\SpecialCharTok}[1]{\textcolor[rgb]{0.25,0.44,0.63}{#1}}
\newcommand{\SpecialStringTok}[1]{\textcolor[rgb]{0.73,0.40,0.53}{#1}}
\newcommand{\StringTok}[1]{\textcolor[rgb]{0.25,0.44,0.63}{#1}}
\newcommand{\VariableTok}[1]{\textcolor[rgb]{0.10,0.09,0.49}{#1}}
\newcommand{\VerbatimStringTok}[1]{\textcolor[rgb]{0.25,0.44,0.63}{#1}}
\newcommand{\WarningTok}[1]{\textcolor[rgb]{0.38,0.63,0.69}{\textbf{\textit{#1}}}}

%% ============================================================
%% DOCUMENT
%% ============================================================
\begin{document}

%% ── Title page ──────────────────────────────────────────────
\frontmatter
\thispagestyle{empty}
\vspace*{2cm}
\begin{center}
  {\fontsize{36}{44}\selectfont\bfseries\sffamily Getting Started with
Docker}\\[0.8em]
  {\Large\sffamily\color{gray} Containers, Images, Networking, Volumes,
Compose and Security}\\[3em]
  {\large Salvatore D'Angelo}\\[0.5em]
  {\normalsize\color{gray} 2025}\\[4em]
  {\small\color{gray} © 2025 Salvatore D'Angelo --- code4projects.com.
All rights reserved.}
\end{center}
\newpage

%% ── Table of contents ───────────────────────────────────────
{
\hypersetup{linkcolor=NavyBlue}
\setcounter{tocdepth}{2}
\tableofcontents
}
\newpage

%% ── Body ────────────────────────────────────────────────────
\mainmatter
\section{Docker Security Best
Practices}\label{docker-security-best-practices}

\emph{Posted on ****}

\subsection{Introduction}\label{introduction}

This is the seventh and final article of the \textbf{Getting Started
with Docker} series. In the
\href{/Users/sasadangelo/github.com/sasadangelo/code4projects/how-docker-compose-works-2025/}{previous
article} we orchestrated a multi-container application with Docker
Compose. Now we turn our attention to \textbf{security}.

Containers are isolated --- but isolation is not the same as security. A
container running as root, with a writable filesystem and full Linux
capabilities, is a significant risk if it is ever compromised. The good
news is that Docker provides several mechanisms to reduce the attack
surface, and most of them require only a few lines of configuration.

\subsection{1. Never Run as Root}\label{never-run-as-root}

By default, processes inside a container run as \textbf{root} (UID 0).
This means that if an attacker exploits a vulnerability in your
application, they have root access inside the container --- and
depending on how the container is configured, they may be able to escape
to the host.

The fix is simple: add a \texttt{USER} instruction to your Dockerfile.

\begin{Shaded}
\begin{Highlighting}[]
\KeywordTok{FROM}\NormalTok{ nginx:alpine}

\KeywordTok{COPY}\NormalTok{ html /usr/share/nginx/html}

\CommentTok{\# Switch to the non{-}root user that the nginx image already provides}
\KeywordTok{USER}\NormalTok{ nginx}

\KeywordTok{EXPOSE}\NormalTok{ 80}
\KeywordTok{CMD}\NormalTok{ [}\StringTok{"nginx"}\NormalTok{, }\StringTok{"{-}g"}\NormalTok{, }\StringTok{"daemon off;"}\NormalTok{]}
\end{Highlighting}
\end{Shaded}

The official \texttt{nginx:alpine} image already includes a non-root
\texttt{nginx} user. For your own images, create a dedicated user:

\begin{Shaded}
\begin{Highlighting}[]
\KeywordTok{FROM}\NormalTok{ alpine:3.20}

\KeywordTok{RUN} \ExtensionTok{addgroup} \AttributeTok{{-}S}\NormalTok{ appgroup }\KeywordTok{\&\&} \ExtensionTok{adduser} \AttributeTok{{-}S}\NormalTok{ appuser }\AttributeTok{{-}G}\NormalTok{ appgroup}

\KeywordTok{COPY} \OperatorTok{{-}{-}chown=appuser:appgroup}\NormalTok{ app /app}

\KeywordTok{USER}\NormalTok{ appuser}

\KeywordTok{CMD}\NormalTok{ [}\StringTok{"/app/start.sh"}\NormalTok{]}
\end{Highlighting}
\end{Shaded}

The \texttt{-\/-chown} flag in \texttt{COPY} ensures the files are owned
by the correct user from the start.

You can also enforce non-root at runtime without changing the
Dockerfile:

\begin{Shaded}
\begin{Highlighting}[]
\NormalTok{docker run {-}{-}user 1001:1001 nginx:alpine}
\end{Highlighting}
\end{Shaded}

\subsection{2. Use a Read-Only
Filesystem}\label{use-a-read-only-filesystem}

If your application does not need to write to its own filesystem, make
it read-only:

\begin{Shaded}
\begin{Highlighting}[]
\NormalTok{docker run {-}{-}read{-}only {-}p 8080:80 nginx:alpine}
\end{Highlighting}
\end{Shaded}

If the application needs to write to specific paths (temp files, logs),
mount those paths as \texttt{tmpfs}:

\begin{Shaded}
\begin{Highlighting}[]
\NormalTok{docker run {-}{-}read{-}only \textbackslash{}}
\NormalTok{  {-}{-}tmpfs /tmp \textbackslash{}}
\NormalTok{  {-}{-}tmpfs /var/cache/nginx \textbackslash{}}
\NormalTok{  {-}{-}tmpfs /var/run \textbackslash{}}
\NormalTok{  {-}p 8080:80 \textbackslash{}}
\NormalTok{  nginx:alpine}
\end{Highlighting}
\end{Shaded}

A read-only filesystem means that even if an attacker gains code
execution inside the container, they cannot modify the application
binaries or install tools.

In Docker Compose:

\begin{Shaded}
\begin{Highlighting}[]
\FunctionTok{services}\KeywordTok{:}
\AttributeTok{  }\FunctionTok{frontend}\KeywordTok{:}
\AttributeTok{    }\FunctionTok{image}\KeywordTok{:}\AttributeTok{ nginx:alpine}
\AttributeTok{    }\FunctionTok{read\_only}\KeywordTok{:}\AttributeTok{ }\CharTok{true}
\AttributeTok{    }\FunctionTok{tmpfs}\KeywordTok{:}
\AttributeTok{      }\KeywordTok{{-}}\AttributeTok{ /tmp}
\AttributeTok{      }\KeywordTok{{-}}\AttributeTok{ /var/cache/nginx}
\AttributeTok{      }\KeywordTok{{-}}\AttributeTok{ /var/run}
\end{Highlighting}
\end{Shaded}

\subsection{3. Drop Linux Capabilities}\label{drop-linux-capabilities}

Linux capabilities are fine-grained divisions of root privileges. By
default, Docker grants a container a subset of them. You can drop all
capabilities and add back only what your application actually needs:

\begin{Shaded}
\begin{Highlighting}[]
\NormalTok{docker run {-}{-}cap{-}drop=ALL {-}{-}cap{-}add=NET\_BIND\_SERVICE nginx:alpine}
\end{Highlighting}
\end{Shaded}

\texttt{NET\_BIND\_SERVICE} allows binding to ports below 1024 (like
port 80). Most application containers do not need any other capability.

In Docker Compose:

\begin{Shaded}
\begin{Highlighting}[]
\FunctionTok{services}\KeywordTok{:}
\AttributeTok{  }\FunctionTok{frontend}\KeywordTok{:}
\AttributeTok{    }\FunctionTok{image}\KeywordTok{:}\AttributeTok{ nginx:alpine}
\AttributeTok{    }\FunctionTok{cap\_drop}\KeywordTok{:}
\AttributeTok{      }\KeywordTok{{-}}\AttributeTok{ ALL}
\AttributeTok{    }\FunctionTok{cap\_add}\KeywordTok{:}
\AttributeTok{      }\KeywordTok{{-}}\AttributeTok{ NET\_BIND\_SERVICE}
\end{Highlighting}
\end{Shaded}

\subsection{4. Never Put Secrets in Environment Variables (or
Images)}\label{never-put-secrets-in-environment-variables-or-images}

A common mistake is passing passwords or API keys through environment
variables in \texttt{docker-compose.yml} or baking them into the image
with \texttt{ENV}. Both approaches expose secrets in
\texttt{docker\ inspect} output and in image layers.

\textbf{What to do instead}:

\begin{itemize}
\tightlist
\item
  Use \textbf{Docker secrets} (available in Swarm mode) which mount
  secrets as files in \texttt{/run/secrets/}
\item
  Use a \textbf{secrets manager} (HashiCorp Vault, AWS Secrets Manager,
  etc.)
\item
  As a minimum, use a \texttt{.env} file that is \textbf{never committed
  to version control} and reference it in Compose
\end{itemize}

\begin{Shaded}
\begin{Highlighting}[]
\CommentTok{\# docker{-}compose.yml}
\FunctionTok{services}\KeywordTok{:}
\AttributeTok{  }\FunctionTok{app}\KeywordTok{:}
\AttributeTok{    }\FunctionTok{image}\KeywordTok{:}\AttributeTok{ my{-}app}
\AttributeTok{    }\FunctionTok{env\_file}\KeywordTok{:}
\AttributeTok{      }\KeywordTok{{-}}\AttributeTok{ .env.secrets}\CommentTok{   \# never committed to git}
\end{Highlighting}
\end{Shaded}

Add to \texttt{.gitignore}:

\begin{verbatim}
.env
.env.secrets
*.env
\end{verbatim}

\subsection{5. Use Minimal Base Images}\label{use-minimal-base-images}

Every package in your base image is a potential attack surface. Prefer
minimal base images:

\begin{itemize}
\tightlist
\item
  \texttt{alpine} --- \textasciitilde5MB, very popular, uses musl libc
\item
  \texttt{distroless} (Google) --- no shell, no package manager, only
  the runtime and your app
\item
  \texttt{scratch} --- completely empty, for statically compiled
  binaries (Go, Rust)
\end{itemize}

For our Nginx example we already use \texttt{nginx:alpine}. Going even
smaller would mean building a custom Nginx binary, which is rarely
necessary.

The rule: \textbf{if it is not needed, it should not be in the image}.

\subsection{6. Scan Images for Known
Vulnerabilities}\label{scan-images-for-known-vulnerabilities}

Even minimal images can contain packages with known CVEs. Docker
provides a built-in scanner called \textbf{Docker Scout}:

\begin{Shaded}
\begin{Highlighting}[]
\NormalTok{\# Quick vulnerability overview}
\NormalTok{docker scout quickview nginx:alpine}

\NormalTok{\# Full CVE list}
\NormalTok{docker scout cves nginx:alpine}

\NormalTok{\# Recommendations to fix vulnerabilities}
\NormalTok{docker scout recommendations nginx:alpine}
\end{Highlighting}
\end{Shaded}

Docker Scout is available in Docker Desktop and as a CLI plugin. It
integrates with Docker Hub and can be added to CI/CD pipelines to fail a
build when high-severity vulnerabilities are detected.

An alternative open-source scanner is \textbf{Trivy}:

\begin{Shaded}
\begin{Highlighting}[]
\NormalTok{trivy image nginx:alpine}
\end{Highlighting}
\end{Shaded}

Make image scanning part of your build pipeline, not an afterthought.

\subsection{7. Keep Images Up to Date}\label{keep-images-up-to-date}

A scanned image with no vulnerabilities today may have CVEs tomorrow.
Pin your base image to a specific digest rather than a floating tag, and
automate rebuild checks:

\begin{Shaded}
\begin{Highlighting}[]
\CommentTok{\# Pin to a specific digest (reproducible, auditable)}
\KeywordTok{FROM}\NormalTok{ nginx:alpine@sha256:...}
\end{Highlighting}
\end{Shaded}

Many teams use tools like \textbf{Renovate} or \textbf{Dependabot} to
automatically open PRs when a newer base image is available.

\subsection{Putting It All Together}\label{putting-it-all-together}

Here is a Dockerfile for our Nginx image that applies all the practices
above:

\begin{Shaded}
\begin{Highlighting}[]
\KeywordTok{FROM}\NormalTok{ nginx:alpine}

\CommentTok{\# Copy content}
\KeywordTok{COPY} \OperatorTok{{-}{-}chown=nginx:nginx}\NormalTok{ html /usr/share/nginx/html}

\CommentTok{\# Drop to non{-}root}
\KeywordTok{USER}\NormalTok{ nginx}

\KeywordTok{EXPOSE}\NormalTok{ 80}

\KeywordTok{CMD}\NormalTok{ [}\StringTok{"nginx"}\NormalTok{, }\StringTok{"{-}g"}\NormalTok{, }\StringTok{"daemon off;"}\NormalTok{]}
\end{Highlighting}
\end{Shaded}

And the corresponding \texttt{docker-compose.yml}:

\begin{Shaded}
\begin{Highlighting}[]
\FunctionTok{services}\KeywordTok{:}
\AttributeTok{  }\FunctionTok{frontend}\KeywordTok{:}
\AttributeTok{    }\FunctionTok{build}\KeywordTok{:}\AttributeTok{ .}
\AttributeTok{    }\FunctionTok{image}\KeywordTok{:}\AttributeTok{ my{-}nginx:secure}
\AttributeTok{    }\FunctionTok{container\_name}\KeywordTok{:}\AttributeTok{ frontend}
\AttributeTok{    }\FunctionTok{ports}\KeywordTok{:}
\AttributeTok{      }\KeywordTok{{-}}\AttributeTok{ }\StringTok{"8080:80"}
\AttributeTok{    }\FunctionTok{read\_only}\KeywordTok{:}\AttributeTok{ }\CharTok{true}
\AttributeTok{    }\FunctionTok{tmpfs}\KeywordTok{:}
\AttributeTok{      }\KeywordTok{{-}}\AttributeTok{ /tmp}
\AttributeTok{      }\KeywordTok{{-}}\AttributeTok{ /var/cache/nginx}
\AttributeTok{      }\KeywordTok{{-}}\AttributeTok{ /var/run}
\AttributeTok{    }\FunctionTok{cap\_drop}\KeywordTok{:}
\AttributeTok{      }\KeywordTok{{-}}\AttributeTok{ ALL}
\AttributeTok{    }\FunctionTok{cap\_add}\KeywordTok{:}
\AttributeTok{      }\KeywordTok{{-}}\AttributeTok{ NET\_BIND\_SERVICE}
\AttributeTok{    }\FunctionTok{networks}\KeywordTok{:}
\AttributeTok{      }\KeywordTok{{-}}\AttributeTok{ app{-}network}

\FunctionTok{networks}\KeywordTok{:}
\AttributeTok{  }\FunctionTok{app{-}network}\KeywordTok{:}
\AttributeTok{    }\FunctionTok{driver}\KeywordTok{:}\AttributeTok{ bridge}
\end{Highlighting}
\end{Shaded}

\subsection{Conclusion}\label{conclusion}

In this article we covered:

\begin{itemize}
\tightlist
\item
  Running containers as a \textbf{non-root user} with \texttt{USER} in
  the Dockerfile
\item
  Using a \textbf{read-only filesystem} with \texttt{-\/-read-only} and
  \texttt{tmpfs} for writable paths
\item
  \textbf{Dropping Linux capabilities} with \texttt{-\/-cap-drop=ALL}
  and adding back only what is needed
\item
  Keeping \textbf{secrets out of images and environment variables}
\item
  Choosing \textbf{minimal base images} to reduce attack surface
\item
  \textbf{Scanning images} for vulnerabilities with Docker Scout or
  Trivy
\item
  \textbf{Keeping images up to date} with pinned digests and automated
  updates
\end{itemize}

This concludes the \textbf{Getting Started with Docker} series. You now
have a solid foundation to build, network, persist, orchestrate, and
secure containerized applications. The natural next step is
\textbf{Docker Swarm} for multi-host deployments, or \textbf{Kubernetes}
for large-scale container orchestration.

\begin{center}\rule{0.5\linewidth}{0.5pt}\end{center}

If you enjoyed this article, don't forget to \textbf{give it a clap 👏},
\textbf{share it with your friends 🔗}, and \textbf{follow me for more
tips and tutorials on software development 📘}. Your support helps me
create more content like this --- thank you! 🙌

\end{document}
